\appendix
\section{Algorithm summaries and size accounting}
\label{sec:appendix}

The following pseudocode fixes the arithmetic decisions independently of
the language-specific implementation. Fractions remain primitive because
each recurrence is unimodular. The integer accumulator is deliberately
not reduced as a rational number.

\begin{lstlisting}
PrimitiveFraction([a0, ..., ak]):
    (p, q) := (1, 0)
    for a in [ak, ..., a0]:
        (p, q) := (a*p + q, p)
    reject unsupported final infinity if q = 0
    if q < 0: (p, q) := (-p, -q)
    return (p, q)

MontesinosDecision(e, source_blocks):
    exceptional := []
    for block in source_blocks:
        (p, q) := PrimitiveFraction(block)
        (a, beta) := EuclideanDivmod(p, q)
        e := e + a
        if q > 1:
            append (beta, q) to exceptional
    (D, Q) := (e, 1)
    for (beta, alpha) in exceptional:
        (D, Q) := (D*alpha + beta*Q, Q*alpha)
    reject input outside the knot domain if D is even
    if len(exceptional) >= 3:
        return KNOTTED, exceptional-fiber certificate
    if abs(D) = 1:
        return UNKNOT, two-bridge certificate
    return KNOTTED, determinant certificate
\end{lstlisting}

The local obstruction checker uses the unnormalized primitive slopes as
well:
\begin{lstlisting}
ForbiddenLocalSource(source):
    require integer term e = 0
    require at least two finite nonintegral rational summands
    compute primitive (pi, qi), with qi > 1
    require two open strands and no closed component
    if number_of_summands >= 3: return FORBIDDEN
    Q := q1*q2
    F := p1*q2 + p2*q1
    return (F mod Q) not in {1, Q - 1}

FindLiteralWitness(host, source):
    compile and verify connected plane pattern with four ports
    require ForbiddenLocalSource(source)
    for target_root in crossings(host):
        for offset in {0, 2}:
            propagate the unique possible map over internal edges
            reject collisions, inconsistent rotations, or overpair errors
            check every remaining dart and four distinct outgoing edges
            if all source and separating-disk checks pass:
                return source, crossing map, rotations
    return UNKNOWN
\end{lstlisting}

\begin{table}[htbp]
\centering\small
\begin{tabularx}{\textwidth}{L{0.30\textwidth}L{0.25\textwidth}X}
\toprule
Operation & Bound & Input and scope\\
\midrule
Two-twist rational summary and gluing &
\(O(m^2)\) bits of time &
Words of length \(m\); \(O(m)\) bits of arithmetic data\\
\addlinespace
Montesinos production arithmetic &
\(O(B^2)\) bits of time &
Complete compressed source length \(B\), including delimiters\\
\addlinespace
Independent branch verification &
\(O(B^3)\) bits of time &
Valid \(O(B)\)-size certificates; parse purported input as well\\
\addlinespace
Literal PD construction &
Polynomial in \(B+C\) &
Must charge \(C\) expanded crossings; cap checked before allocation\\
\addlinespace
One compiled pattern search &
\(O(nm)\) word operations &
Connected \(m\)-crossing pattern, indexed target already validated\\
\addlinespace
Public finite-catalogue wrapper &
\(O(n\log(n+2)+nM)\) word operations &
Total expanded pattern size \(M\), plus all source/compilation costs\\
\addlinespace
Bounded-description discovery &
\((n+2)^{O(h)}\) bit operations &
Literal admissible sources with at most \(h\) coefficients\\
\bottomrule
\end{tabularx}
\caption{Complexity statements and the input sizes they actually charge.
The abstract matching word bounds have indexed-array realizations.
The Python implementation's table operations use the usual table model.}
\label{tab:complexity}
\end{table}

For the expansion row, even a conservative implementation that revisits
the current explicit graph at each construction step performs polynomially
many operations in \(B+C\). The delivered compiler avoids unnecessary
whole-graph mirroring, but this report makes no finer worst-case compiler
claim than is needed. The production arithmetic bound applies before
any expansion and is the bound relevant to exponentially large twist
coefficients.

\section{Artifact map and reproducibility}

The main article driver is
\code{article/unknot_arithmetic_continuations.tex}. It is assembled from
numbered editable sections by \code{article/build_article.py}; all prose,
theorems, tables, and bibliography are included directly in the combined
driver. The optional plotted figure is used when its file exists, so
the standalone combined source can also compile without the figure assets.
The archive contains both the assets and a portable plotting script.

\begin{table}[htbp]
\centering\small
\begin{tabularx}{\textwidth}{L{0.41\textwidth}X}
\toprule
Path in the archive & Purpose\\
\midrule
\code{article/} & Combined and section TeX, PDF, build script,
supplementary BibTeX, and figure assets\\
\code{fast/} & Runnable updated recognizer, standard-library implementation,
tests, examples, certificates, and preserved license\\
\code{integration/} & Reviewable patch, application instructions, source
manifest, and patch byte-verification manifest\\
\code{experiments/} & Corpus generator, paired benchmark,
independent arithmetic and topology diagnostics, reproduction launcher\\
\code{examples/} & Complete structured and planted plain-PD inputs\\
\code{certificates/} & Complete local source-plus-map certificates\\
\code{results/} & Raw timings, finite verification outputs, full-suite
and patch-application logs\\
\code{RESEARCH_PROVENANCE.md} & Baseline, imported results,
independently derived claims, and limited inspiration attribution\\
\bottomrule
\end{tabularx}
\caption{Files supplied for review and later integration into ProveIt.}
\end{table}

The reproduction launcher accepts an explicit output directory and creates
an isolated replay workspace, leaving the bundled measurement record
unchanged. From any working directory, invoke its absolute path with
\code{--quick} for the added suites and finite diagnostics,
\code{--full-tests} for the full suite, \code{--identity} for all identity
matrices through \(m=10\), or \code{--benchmarks} for a fresh seven-round
timing run. The \code{--all} mode intentionally performs the expensive
reproductions as well. Exact examples are in the top-level README.

The repository patch was created for the pinned baseline, not for an
unseen future revision. Its manifest permits a reviewer to compare the
expected bytes before application. Subsequent repository changes may
require a normal patch merge and another targeted validation; that is
an integration issue rather than a change in any theorem above.

\begin{thebibliography}{99}
\small\raggedright

\bibitem{proveit-baseline}
V. Reshetnikov and contributors.
\emph{ProveIt: Topology/UnknotRecognition}.
Repository source and research reports, commit
\code{86f23490006c880cd01db9bec10761d39c6395c8},
inspected 8 October 2026.
\url{https://github.com/VladimirReshetnikov/ProveIt/tree/86f23490006c880cd01db9bec10761d39c6395c8/Topology/UnknotRecognition}.

\bibitem{kronheimer-mrowka}
P. B. Kronheimer and T. S. Mrowka.
\emph{Khovanov homology is an unknot-detector}.
arXiv:1005.4346, 2010.
\url{https://arxiv.org/abs/1005.4346}.

\bibitem{lackenby-2021}
M. Lackenby.
\emph{Unknot recognition in quasi-polynomial time}.
Oxford lecture slides, March 2021, especially slides 7 and 25--34.
The displayed announced bound is \(2^{O((\log n)^3)}\).
\url{https://people.maths.ox.ac.uk/lackenby/quasipolynomial-talk-oxford-compressed.pdf}.

\bibitem{lackenby-2026}
M. Lackenby.
\emph{Incompressible surfaces, hierarchies and unknot recognition}.
arXiv:2607.23350v1, 25 July 2026, especially Sections 9--10.
\url{https://arxiv.org/abs/2607.23350}.

\bibitem{kauffman-lambropoulou-tangles}
L. H. Kauffman and S. Lambropoulou.
On the classification of rational tangles.
\emph{Advances in Applied Mathematics} 33 (2004), 199--237.
arXiv:math/0311499.
\url{https://arxiv.org/abs/math/0311499}.

\bibitem{kauffman-lambropoulou-knots}
L. H. Kauffman and S. Lambropoulou.
On the classification of rational knots.
\emph{L'Enseignement Math\'ematique} 49 (2003), 357--410.
arXiv:math/0212011.
\url{https://arxiv.org/abs/math/0212011}.

\bibitem{KauffmanLambropoulouHard}
L. H. Kauffman and S. Lambropoulou.
\emph{Hard unknots and collapsing tangles}.
arXiv:math/0601525v5, 26 July 2014; original preprint 2006.
Especially Theorems 4--5 and the rational fraction conventions.
\url{https://arxiv.org/abs/math/0601525}.

\bibitem{KauffmanLambropoulouFractions}
L. H. Kauffman and S. Lambropoulou.
\emph{From tangle fractions to DNA}.
arXiv:math/0311511, 2003, especially Theorem 1 and Section 4.1.
\url{https://arxiv.org/abs/math/0311511}.

\bibitem{nogueira-salgueiro}
J. M. Nogueira and A. Salgueiro.
On embeddings of 2-string tangles into the unknot, the unlink and split links.
\emph{Topology and its Applications} 305 (2022), 107866.
arXiv:2110.15645v1, especially Proposition 4.4 and Theorem 4.9.
\url{https://arxiv.org/abs/2110.15645}.

\bibitem{belousov-pretzel}
Y. Belousov.
\emph{Explicit formulas for the Alexander polynomial of pretzel knots}.
arXiv:2502.10370v4, 9 March 2026.
\emph{Bulletin of the Australian Mathematical Society},
doi:10.1017/S0004972726101439.
Especially Theorem 1(3) and Section 2.3.
\url{https://arxiv.org/abs/2502.10370}.

\bibitem{BarNatanTangles}
D. Bar--Natan.
Khovanov's homology for tangles and cobordisms.
\emph{Geometry \& Topology} 9 (2005), 1443--1499.
arXiv:math/0410495, especially Theorem 2 and Section 7.
\url{https://arxiv.org/abs/math/0410495}.

\bibitem{KhovanovPatterns}
M. Khovanov.
Patterns in knot cohomology I.
\emph{Experimental Mathematics} 12 (2003), no.~3, 365--374.
arXiv:math/0201306, especially Section 3.
\url{https://arxiv.org/abs/math/0201306}.

\bibitem{ShumakovitchTorsion}
A. N. Shumakovitch.
Torsion of the Khovanov homology.
\emph{Fundamenta Mathematicae} 225 (2014), 343--364.
arXiv:math/0405474v2, 19 June 2018, Corollary 3.2.C.
\url{https://arxiv.org/abs/math/0405474}.

\bibitem{OpenAIMathContextSeparation}
OpenAI.
\emph{An exponential two-way deterministic state lower bound for one-way
liveness}, inspected source section on relation products and context
distinguishability. In the \code{openai/math} repository, commit
\code{adc7f1241b42e322a6451854ab7e4b4c146bf78a}, 2026.
Used only as inspiration for the standard context-separation argument;
no headline automata theorem is assumed here.
\url{https://github.com/openai/math/tree/adc7f1241b42e322a6451854ab7e4b4c146bf78a}.
The precise source path and blob hash are in the accompanying provenance.

\end{thebibliography}
\end{document}
